\documentclass[11pt]{amsart}
\usepackage[T1]{fontenc}\usepackage{lmodern}\usepackage{amsmath,amssymb,amsthm}\usepackage{xcolor}
\usepackage[colorlinks=true,linkcolor=blue!55!black,citecolor=blue!55!black,urlcolor=blue!55!black]{hyperref}
\newtheorem{theorem}{Theorem}\newtheorem{lemma}[theorem]{Lemma}
\title{A Certified Prime Factor of \texorpdfstring{$a_{20}$}{a(20)} in the Recurrence \texorpdfstring{$a_{n+1}=a_n^2+2$}{a(n+1)=a(n) squared plus 2}}
\author{Xiaojian Zeng}\address{PBC School of Finance, Tsinghua University, 100190}\email{xiaojianzeng@foxmail.com}
\date{September 2, 2026}
\hypersetup{pdftitle={A Certified Prime Factor of a(20) in the Recurrence a(n+1)=a(n) squared plus 2},pdfauthor={Xiaojian Zeng}}
\begin{document}
\begin{abstract}
Let $a_0=1$ and $a_{n+1}=a_n^2+2$. The OEIS record A102847 reports
$a_8,\ldots,a_{19}$ composite and the primality of $a_{20}$ unknown. We give
the exact certificate $7\,795\,379\mid a_{20}$. The divisor is prime. Together
with an elementary congruence for odd indices, this proves that every $a_n$ with
$5\le n\le21$ is composite. Thus a next prime after $a_4=15131$, if one
exists, has index at least $22$. No assertion is made for $n\ge22$.
\end{abstract}\maketitle
\section{Statement and prior frontier}
The integer sequence A102847 is defined by
\begin{equation}\label{eq:rec}a_0=1,\qquad a_{n+1}=a_n^2+2\quad(n\ge0).\end{equation}
The OEIS record lists prime indices $1,2,4$, reports indices $8$ through $19$
composite, and records the primality of $a_{20}$ as unknown \cite{OEIS102847}.
A Formal Conjectures repository source repeats the definition and poses the
next-prime question as open \cite{Formal102847}.
\begin{theorem}\label{thm:main}
For \eqref{eq:rec}, $a_n$ is composite for every $5\le n\le21$. In particular,
if there is a least prime term after $a_4$, its index is at least $22$.
\end{theorem}
\section{Proof and exact certificate}
\begin{lemma}\label{lem:odd}For every odd $n\ge1$, $3\mid a_n$.\end{lemma}
\begin{proof}Modulo $3$, the recurrence begins $1,0,2,0,2,0,\ldots$.
A zero is sent to $2$, and $2$ is sent to $0$; induction proves the claim.\end{proof}
\begin{proof}[Proof of Theorem \ref{thm:main}]
For $5\le n\le20$, the table gives one prime divisor of $a_n$:
\[
\begin{array}{c|rrrrrrrr}n&5&6&7&8&9&10&11&12\\ \hline
p&3&19&3&41&3&612304337&3&19\\[2pt]
n&13&14&15&16&17&18&19&20\\ \hline
p&3&182094337&3&43&3&19&3&7795379\end{array}
\]
Each entry is checked by iterating \eqref{eq:rec} modulo $p$. For the new
endpoint, the complete residue chain modulo $7\,795\,379$, from $a_0$ through
$a_{20}$, is
\begin{align*}
&1,3,11,123,15131,2881172,1108687,3707872,2390415,906816,\\
&4113285,7318869,5775969,4541074,3143408,4589153,1945093,\\
&920928,2327502,7240778,0.
\end{align*}
Trial division through $\lfloor\sqrt{7\,795\,379}\rfloor=2792$ proves the
new divisor prime. Every divisor in the table is proper: for $n=5$ this is
$3<a_5=228\,947\,163$, while for $6\le n\le20$ it follows from
\[
  p\le 612\,304\,337<a_6=52\,416\,803\,445\,748\,571\le a_n.
\]
Lemma \ref{lem:odd} covers $n=21$.\end{proof}
\section{Reproducibility and limitation}
The accompanying JSON contains the index--divisor pairs. The independent
verifier checks primality by deterministic trial division, recomputes every
modular recurrence, and checks properness from the increasing exact recurrence.
It uses Python 3 exact integer arithmetic and no external primality oracle.
The discovery search tested prime divisors at most $20\,000\,000$ for indices
through $22$; finding none for $a_{22}$ is not evidence of primality and is not
used in the theorem.

This is a certified finite extension only. It neither produces a prime term
after $a_4$ nor proves that none exists. The literature search was limited to
the recorded sources and queries through September 2, 2026.
\section*{Data and computational disclosure}
The source archive includes the exact certificate, verifier, and raw logs.
Exact integer computation was used to discover and verify the
factor and prepare the reproducibility files. The theorem can be checked from
the displayed recurrence and finite integer certificate.
\bibliographystyle{amsplain}\bibliography{references}
\end{document}
